\documentclass[10pt]{beamer}
\usetheme{metropolis}
\usepackage{graphicx}
\usepackage{listings}
\usepackage{booktabs}
\usepackage{xcolor}
\usepackage{hyperref}

\definecolor{codebg}{HTML}{F5F5F5}
\definecolor{codegreen}{HTML}{2E7D32}
\definecolor{codegray}{HTML}{757575}
\definecolor{codeblue}{HTML}{1565C0}

\lstdefinestyle{rstyle}{
  language=R, backgroundcolor=\color{codebg}, basicstyle=\ttfamily\scriptsize,
  keywordstyle=\color{codeblue}\bfseries, stringstyle=\color{codegreen},
  commentstyle=\color{codegray}\itshape, breaklines=true, frame=single,
  rulecolor=\color{codegray!30}, showstringspaces=false, tabsize=2,
  morekeywords={library,ggplot,aes,geom_point,geom_line,geom_col,geom_bar,
    geom_smooth,geom_histogram,geom_hex,
    scale_x_continuous,scale_x_log10,scale_x_discrete,scale_x_date,
    scale_y_continuous,scale_y_log10,
    scale_color_manual,scale_color_brewer,scale_color_viridis_c,
    scale_color_gradient2,scale_fill_viridis_c,scale_fill_brewer,
    scale_size,coord_cartesian,coord_flip,coord_polar,coord_fixed,
    theme_minimal,theme,element_text,element_blank,labs,
    date_breaks,date_labels,comma,dollar,percent,label_comma}
}
\lstdefinestyle{pythonstyle}{
  language=Python, backgroundcolor=\color{codebg}, basicstyle=\ttfamily\scriptsize,
  keywordstyle=\color{codeblue}\bfseries, stringstyle=\color{codegreen},
  commentstyle=\color{codegray}\itshape, breaklines=true, frame=single,
  rulecolor=\color{codegray!30}, showstringspaces=false, tabsize=4,
  morekeywords={import,from,as,pd,plt,sns,np,fig,ax,set_xscale,set_yscale,
    set_xlim,set_ylim,set_aspect,FuncFormatter,DateFormatter,YearLocator,
    TwoSlopeNorm}
}

\title{Scales \& Coordinate Systems}
\subtitle{Workshop 2 --- Module 4}
\author{Prof.Asoc.\ Endri Raco}
\institute{Data Visualization for Data Scientists \\ UNYT Tirana}
\date{2025/26}

\begin{document}

\begin{frame}
  \titlepage
\end{frame}

\begin{frame}{Module Roadmap}
  \begin{enumerate}
    \item Anatomy of a scale: domain $\rightarrow$ transform $\rightarrow$ visual range
    \item Position scales: continuous, log, discrete, date
    \item Colour scales: manual, Brewer, viridis, gradient2
    \item Axis label formatting: comma, dollar, percent
    \item Coordinate systems: Cartesian, flip, polar, fixed
    \item Zooming safely: \texttt{coord\_cartesian} vs scale \texttt{limits}
  \end{enumerate}
  \begin{exampleblock}{Learning Outcome}
    You will control how every data value maps to a visual property ---
    position, colour, size --- and how the coordinate system frames
    the chart.
  \end{exampleblock}
\end{frame}

\section{Anatomy of a Scale}

\begin{frame}{Domain $\rightarrow$ Transform $\rightarrow$ Visual Range}
  \begin{center}
    \includegraphics[width=0.92\textwidth]{figures_w02m04/scale_anatomy}
  \end{center}
  \begin{exampleblock}{Data Insight}
    A scale answers three questions: (1) What are the input values?
    (domain), (2) Should they be transformed? (log, sqrt, reverse),
    (3) What visual output do they produce? (pixel position, hex colour,
    point size).
  \end{exampleblock}
\end{frame}

\begin{frame}{Scale Functions Cheat Sheet}
  \begin{center}
    \includegraphics[width=0.88\textwidth]{figures_w02m04/scale_types}
  \end{center}
\end{frame}

\section{Position Scales}

\begin{frame}{Linear vs Log Scale}
  \begin{center}
    \includegraphics[width=0.92\textwidth]{figures_w02m04/linear_vs_log}
  \end{center}
  \begin{alertblock}{Pro-Tip}
    Use \texttt{scale\_x\_log10()} (R) or \texttt{ax.set\_xscale("log")}
    (Python) whenever data spans more than two orders of magnitude.
    The log transform converts multiplicative relationships to linear
    ones, making patterns visible.
  \end{alertblock}
\end{frame}

\begin{frame}[fragile]{Log Scale in R vs Python}
  \begin{columns}[T]
    \begin{column}{0.48\textwidth}
      \textbf{R (ggplot2)}
\begin{lstlisting}[style=rstyle]
ggplot(apps, aes(x = Reviews,
                  y = Rating)) +
  geom_point(alpha = 0.3,
             size = 0.8) +

  # Log-transform x axis
  scale_x_log10(
    labels = scales::comma) +

  # Limit y without filtering
  coord_cartesian(
    ylim = c(1, 5)) +

  theme_minimal() +
  labs(x = "Reviews (log scale)")
\end{lstlisting}
    \end{column}
    \begin{column}{0.48\textwidth}
      \textbf{Python (matplotlib)}
\begin{lstlisting}[style=pythonstyle]
fig, ax = plt.subplots()

ax.scatter(reviews, rating,
    s=8, alpha=0.3)

# Log-transform x axis
ax.set_xscale("log")

# Zoom without filtering
ax.set_ylim(1, 5)

# Format tick labels
from matplotlib.ticker import (
    FuncFormatter)
ax.xaxis.set_major_formatter(
    FuncFormatter(lambda x, p:
        f"{x:,.0f}"))

ax.set_xlabel("Reviews (log)")
\end{lstlisting}
    \end{column}
  \end{columns}
\end{frame}

\begin{frame}{Date Axis Formatting}
  \begin{center}
    \includegraphics[width=0.92\textwidth]{figures_w02m04/date_scale}
  \end{center}
\end{frame}

\begin{frame}[fragile]{Date Scales in R vs Python}
  \begin{columns}[T]
    \begin{column}{0.48\textwidth}
      \textbf{R}
\begin{lstlisting}[style=rstyle]
ggplot(df, aes(x = date,
               y = value)) +
  geom_line() +
  scale_x_date(
    date_breaks = "6 months",
    date_labels = "%b %Y") +
  # Produces: "Jan 2023"
  theme_minimal()
\end{lstlisting}
    \end{column}
    \begin{column}{0.48\textwidth}
      \textbf{Python}
\begin{lstlisting}[style=pythonstyle]
import matplotlib.dates as mdates

ax.plot(dates, values)

ax.xaxis.set_major_locator(
    mdates.MonthLocator(
        bymonth=[1, 7]))
ax.xaxis.set_major_formatter(
    mdates.DateFormatter("%b %Y"))

fig.autofmt_xdate()
\end{lstlisting}
    \end{column}
  \end{columns}
\end{frame}

\section{Colour Scales}

\begin{frame}{Axis Label Formatting}
  \begin{center}
    \includegraphics[width=0.92\textwidth]{figures_w02m04/axis_labels}
  \end{center}
  \begin{alertblock}{Pro-Tip}
    In R: \texttt{scale\_y\_continuous(labels = scales::comma)} or
    \texttt{scales::dollar} or \texttt{scales::percent}. In Python:
    \texttt{ax.yaxis.set\_major\_formatter(FuncFormatter(...))} or
    use \texttt{matplotlib.ticker.StrMethodFormatter}.
  \end{alertblock}
\end{frame}

\begin{frame}[fragile]{Colour Scale Selection}
  \centering
  \begin{tabular}{lll}
    \toprule
    \textbf{Data Type} & \textbf{R Scale} & \textbf{Python Equivalent} \\
    \midrule
    Categorical ($\leq$7) & \texttt{scale\_color\_brewer("Set2")} & \texttt{palette="Set2"} \\
    Categorical (custom) & \texttt{scale\_color\_manual(values)} & loop + palette dict \\
    Sequential cont. & \texttt{scale\_color\_viridis\_c()} & \texttt{cmap="viridis"} \\
    Diverging cont. & \texttt{scale\_color\_gradient2()} & \texttt{TwoSlopeNorm} \\
    Binary & \texttt{scale\_color\_manual(c(a,b))} & 2-colour list \\
    \bottomrule
  \end{tabular}
  \vspace{4pt}

  {\small Always match scale type to variable type. Using a sequential
  scale for categorical data implies a false ordering; using a
  qualitative scale for continuous data creates false boundaries.}
\end{frame}

\begin{frame}{Diverging vs Sequential}
  \begin{center}
    \includegraphics[width=0.88\textwidth]{figures_w02m04/diverging}
  \end{center}
  \begin{exampleblock}{Data Insight}
    Use a diverging scale (\texttt{RdBu\_r}, \texttt{gradient2}) when
    the data has a meaningful midpoint (zero, average, break-even).
    The two hues encode direction; the midpoint anchors neutral.
    Use sequential (\texttt{viridis}) when only magnitude matters.
  \end{exampleblock}
\end{frame}

\section{Coordinate Systems}

\begin{frame}{Four Coordinate Systems}
  \begin{center}
    \includegraphics[width=0.95\textwidth]{figures_w02m04/coord_systems}
  \end{center}
\end{frame}

\begin{frame}[fragile]{Coordinate Systems in Code}
\begin{lstlisting}[style=rstyle]
# Default: Cartesian
p + coord_cartesian()

# Flip x and y (horizontal bars)
p + coord_flip()

# Polar (bar -> coxcomb/pie)
p + coord_polar(theta = "y")   # pie: maps y to angle
p + coord_polar(theta = "x")   # coxcomb: maps x to angle

# Equal aspect ratio (maps, scatter comparison)
p + coord_fixed(ratio = 1)

# Python equivalents:
# coord_flip    -> ax.barh() instead of ax.bar()
# coord_polar   -> fig.add_subplot(projection="polar")
# coord_fixed   -> ax.set_aspect("equal")
\end{lstlisting}
\end{frame}

\section{Zooming Safely}

\begin{frame}{scale limits vs coord\_cartesian}
  \begin{center}
    \includegraphics[width=0.95\textwidth]{figures_w02m04/zoom}
  \end{center}
\end{frame}

\begin{frame}[fragile]{The Zoom Rule}
  \begin{columns}[T]
    \begin{column}{0.48\textwidth}
      \textbf{DANGEROUS: scale limits}
\begin{lstlisting}[style=rstyle]
# Filters data outside range!
ggplot(df, aes(x, y)) +
  geom_point() +
  geom_smooth(method = "lm") +
  scale_x_continuous(
    limits = c(25, 75))
# geom_smooth recalculated on
# subset -> different slope!
\end{lstlisting}
      {\small Data outside limits is \textbf{removed}
      before statistics are computed. The regression
      line changes.}
    \end{column}
    \begin{column}{0.48\textwidth}
      \textbf{SAFE: coord\_cartesian}
\begin{lstlisting}[style=rstyle]
# Zooms without filtering
ggplot(df, aes(x, y)) +
  geom_point() +
  geom_smooth(method = "lm") +
  coord_cartesian(
    xlim = c(25, 75))
# geom_smooth computed on ALL
# data, then clipped to window
\end{lstlisting}
      {\small Data is only \textbf{clipped} at the
      viewport edge. Statistics remain unchanged.}
    \end{column}
  \end{columns}

  \vspace{6pt}
  \begin{alertblock}{Pro-Tip}
    \textbf{Always} use \texttt{coord\_cartesian()} for zooming. Use
    scale \texttt{limits} only when you intentionally want to
    \textbf{exclude} data from computation.
    In Python: \texttt{ax.set\_xlim()} is equivalent to
    \texttt{coord\_cartesian} --- it clips, not filters.
  \end{alertblock}
\end{frame}

\section{Complete Examples}

\begin{frame}[fragile]{R: Combining Scale + Coord}
  \begin{columns}[T]
    \begin{column}{0.52\textwidth}
\begin{lstlisting}[style=rstyle]
ggplot(apps_clean,
  aes(x = Reviews, y = Rating,
      color = Type)) +
  geom_point(alpha = 0.3,
             size = 0.8) +

  # Position scale: log
  scale_x_log10(
    labels = scales::comma) +

  # Colour scale: manual
  scale_color_manual(
    values = c(Free = "#1565C0",
               Paid = "#E53935")) +

  # Zoom safely
  coord_cartesian(
    ylim = c(1, 5)) +

  theme_minimal() +
  labs(title = "Reviews vs Rating",
       x = "Reviews (log)",
       y = "Rating")
\end{lstlisting}
    \end{column}
    \begin{column}{0.45\textwidth}
      \includegraphics[width=\textwidth]{figures_w02m04/r_scale}

      \vspace{4pt}
      {\small Three scale operations:
      (1) x-axis log-transformed,
      (2) colour manually assigned,
      (3) y-axis zoomed via coord.}
    \end{column}
  \end{columns}
\end{frame}

\begin{frame}[fragile]{Python: Same Three Operations}
  \begin{columns}[T]
    \begin{column}{0.52\textwidth}
\begin{lstlisting}[style=pythonstyle]
fig, ax = plt.subplots(figsize=(6,5))

palette = {"Free": "#1565C0",
           "Paid": "#E53935"}
for t, c in palette.items():
    mask = df["Type"] == t
    ax.scatter(
        df.loc[mask, "Reviews"],
        df.loc[mask, "Rating"],
        s=8, c=c, alpha=0.3,
        edgecolors="none", label=t)

# Position scale: log
ax.set_xscale("log")

# Colour: handled by palette dict

# Zoom: set_ylim clips, not filters
ax.set_ylim(1, 5)

ax.legend(fontsize=7)
ax.set_xlabel("Reviews (log)")
ax.set_ylabel("Rating")
ax.spines["top"].set_visible(False)
ax.spines["right"].set_visible(False)
\end{lstlisting}
    \end{column}
    \begin{column}{0.45\textwidth}
      \includegraphics[width=\textwidth]{figures_w02m04/py_scale}

      \vspace{4pt}
      {\small matplotlib equivalents:
      \texttt{set\_xscale("log")} = log,
      palette dict = manual colour,
      \texttt{set\_ylim()} = safe zoom.}
    \end{column}
  \end{columns}
\end{frame}

\section{Synthesis}

\begin{frame}{Scale Selection Decision Tree}
  \centering
  \begin{tabular}{lll}
    \toprule
    \textbf{Question} & \textbf{Answer} & \textbf{Scale} \\
    \midrule
    Spans $>$2 orders?     & Yes  & \texttt{scale\_x\_log10()} \\
    Has a meaningful zero? & Yes  & \texttt{scale\_color\_gradient2()} \\
    Categorical $\leq$7?   & Yes  & \texttt{scale\_color\_brewer()} \\
    Dates on x?            & Yes  & \texttt{scale\_x\_date()} \\
    Need horizontal bars?  & Yes  & \texttt{coord\_flip()} \\
    Need equal aspect?     & Yes  & \texttt{coord\_fixed(ratio=1)} \\
    Want to zoom?          & Yes  & \texttt{coord\_cartesian(xlim)} \\
    \bottomrule
  \end{tabular}
\end{frame}

\begin{frame}{Key Takeaways}
  \begin{enumerate}
    \item A \textbf{scale} maps data values (domain) through an optional
          transform to visual properties (range).
    \item \textbf{Log scales} reveal patterns in skewed data; use when
          values span $>$2 orders of magnitude.
    \item \textbf{Diverging} colour scales need a meaningful midpoint;
          \textbf{sequential} scales encode magnitude only.
    \item Format axis labels with \texttt{scales::comma/dollar/percent}
          (R) or \texttt{FuncFormatter} (Python).
    \item \textbf{coord\_cartesian} zooms safely (clips); scale
          \texttt{limits} zoom dangerously (filters data).
    \item In Python, \texttt{ax.set\_xlim()} is a safe zoom (like
          \texttt{coord\_cartesian}), not a filter.
  \end{enumerate}
\end{frame}

\begin{frame}{Essential Readings}
  \begin{itemize}
    \item Wickham, H. (2016). \emph{ggplot2}, Chapters 6--8
          (Scales, Colour, Coordinate Systems).
    \item Wilke, C.~O. (2019). \emph{Fundamentals of Data Visualization},
          Chapters 3, 7, 8 (Coordinate Systems, Colour Scales).
    \item Wickham, H. \& Grolemund, G. (2023). \emph{R for Data Science},
          2nd ed., Chapter 11 (Communication: Scales).
  \end{itemize}
  \vspace{6pt}
  \textbf{Next Module}: Faceting \& Small Multiples (W02-M05)
\end{frame}

\end{document}
